\section{导论：两百种语言不是一个模型数字}

这堂课讨论的不是“把 Transformer 再放大一点”，而是一个更难的系统问题：当可用数据、书写传统、语言标准、评测资源与真实需求都极不均衡时，怎样把机器翻译从少数高资源语言扩展到两百种语言，同时保住质量、安全性与可用性。NLLB（No Language Left Behind）把这个目标拆成一条完整链路：先访谈使用者，再建设评测集与种子数据，随后挖掘语料、训练多语模型、抑制过拟合，最后用自动指标、人工评测和 toxicity 检查共同验收。

\lecturefigure{slide-01-title.jpg}{No Language Left Behind：从语言覆盖走向人本机器翻译}{Stanford 课堂视频 00:00:54}

\begin{knowledgebox}{读图：标题里的三个关键词}
\term{No Language Left Behind} 不是保证每种语言都已得到同等质量，而是研究方向；\term{Scaling} 不只指参数量，还包括语言数量、翻译方向、数据管线和评测能力；\term{Human-Centered} 则要求需求定义、数据标准、人工验收和安全边界都由真实使用者参与。整堂课的结构正是围绕这三层展开。
\end{knowledgebox}

\teachervoice{Fan 开场说，她原本也可以讲 Llama 2，但 NLLB 与课程的生成式 AI 主题同样相关，而且这个问题对她有个人意义：英语是她的第三语言。这个动机提醒我们，多语系统不是抽象 benchmark，而是决定谁能被技术服务。}

\begin{importantbox}{本讲的系统命题}
多语机器翻译的质量不是单一模型函数，而是一个耦合系统：
\[
Q = F(H,E,S,M,V),
\]
其中 $H$ 表示 human needs 与参与机制，$E$ 表示 evaluation infrastructure，$S$ 表示 source data 与 mining pipeline，$M$ 表示 model/training，$V$ 表示 validation 与 safety。任一环节缺失，增加参数或数据量都可能只扩大已有偏差。
\end{importantbox}

\subsection{阅读路线}

本讲先解释为什么“支持多少语言”不能单独代表进步；随后沿课堂顺序讨论访谈、FLORES-200、NLLB-Seed、WikiMatrix/LASER3/Stopes、数据配方、MoE 与 curriculum learning；最后把自动指标、人类判断、toxicity 和课堂 Q\&A 统一成一套可审计的评估框架。

\subsection{本章小结}

NLLB 的真正研究对象是“语言覆盖系统”，不是孤立模型。后文每一个数据集、架构和指标，都要放回需求、证据与安全的完整闭环中理解。

